\documentclass[11pt,a4paper]{article}
% Common preamble for the hint sheets (\input from each file)
\usepackage[utf8]{inputenc}
\usepackage[T1]{fontenc}
\usepackage{lmodern}
\usepackage{amsmath,amssymb,amsthm,mathtools}
\usepackage[a4paper,margin=2.3cm]{geometry}
\usepackage{enumitem}
\usepackage{xcolor}
\usepackage[hidelinks]{hyperref}
\renewcommand{\leq}{\leqslant}
\renewcommand{\le}{\leqslant}
\renewcommand{\geq}{\geqslant}
\renewcommand{\ge}{\geqslant}
\setlength{\parindent}{0pt}
\setlength{\parskip}{0.35em}
% \hintstitle{Sheet title}{Class work / Homework}
\newcommand{\hintstitle}[2]{\begin{center}{\LARGE\bfseries Hints: #1\par}\vspace{0.5em}{\large #2\par}\end{center}
\noindent\textit{The hints to each problem are cumulative. Read one, try again, and open the next only if you are still stuck.}\par\medskip\hrule\medskip}
\newcommand{\problem}[1]{\par\bigskip{\large\bfseries Problem #1}\par\nopagebreak\smallskip}
\newcommand{\subpart}[1]{\par\smallskip\textbf{(#1)}\par\nopagebreak}
\newcommand{\hint}[1]{\par\smallskip\noindent\textcolor{black!70}{\textbf{Hint #1.}}\enspace}
\newcommand{\note}[1]{\par\smallskip\noindent\textit{Note.} #1\par}
\providecommand{\lt}{<}
\providecommand{\gt}{>}
\providecommand{\tightlist}{\setlength{\itemsep}{0pt}\setlength{\parskip}{0pt}}
\emergencystretch=3em

\begin{document}
\hintstitle{Methods of proof -- 1}{Class work}

\problem{1}
\hint{1} For the pigeonhole: if every box held at most $\lceil n/k\rceil-1$ objects, the total would be at most $k(\lceil n/k\rceil-1)<n$. For the approximation, look at the $N+1$ numbers $\{0\cdot\alpha\},\{\alpha\},\dots,\{N\alpha\}$ and the $N$ boxes $[j/N,(j+1)/N)$, $j=0,\dots,N-1$.
\hint{2} Two of them, $\{a\alpha\}$ and $\{b\alpha\}$ with $a<b$, share a box. With $q=b-a\le N$ this gives an integer $p$ with $|q\alpha-p|<1/N$, so $\|q\alpha\|<1/N$, and then $|\alpha-p/q|<1/(qN)\le 1/q^2$. To get infinitely many fractions: since $\alpha$ is irrational, $|\alpha-p/q|>0$, so choose $N$ with $1/N$ smaller than this error and get a new fraction.
\hint{3} Density: given an interval of length $\varepsilon$ in $[0,1]$, take $N>1/\varepsilon$ and $q$ with $\delta=\|q\alpha\|<1/N$, $\delta>0$. The points $\{mq\alpha\}$, $m=1,2,\dots$, then move along the circle $\mathbb R/\mathbb Z$ in steps of exactly $\delta$ (forwards if $\{q\alpha\}=\delta$, backwards if $\{q\alpha\}=1-\delta$). Steps shorter than $\varepsilon$ cannot jump over the interval.

\problem{2}
\hint{1} Blichfeldt: integrate $\sum_i \mathbf 1_{A_i}$ over the ambient set of measure $\mu$. If every point lay in at most $k$ of the sets, the integral would be at most $k\mu$. In the lattice version, translate a set $S\subseteq\mathbb R^n$ of volume $>1$ into the unit cube $[0,1)^n$ piece by piece (one piece for each integer translate). The pieces have total volume $>1$, so two points $x\neq y$ of $S$ have $x-y\in\mathbb Z^n$.
\hint{2} Minkowski: apply this to $S=\tfrac12 K$, which has volume $>1$ when $\operatorname{vol}K>2^n$. Then $x-y=\tfrac12(2x)+\tfrac12(-2y)\in K$ by symmetry and convexity, and $x-y$ is a non-zero lattice point.
\hint{3} Steinhaus: we may shrink $A$ to a bounded set of positive measure inside an interval $I$. For $|t|<\delta$, the sets $A$ and $A+t$ both lie in an interval of length $|I|+\delta$. If $2m(A)>|I|+\delta$ they must overlap, which gives $t\in A-A$. To arrange $2m(A)>|I|+\delta$, first replace $A$ by $A\cap J$ for an interval $J$ where $A$ occupies more than half of $J$. Such a $J$ exists by regularity: cover $A$ by an open set $U$ with $m(U)<2m(A)$ and look at its component intervals.
\hint{4} For the additive $f$: if $|f|\le M$ on $A$, then $|f(a-a')|\le 2M$, so $f$ is bounded on an interval $(-\delta,\delta)$. Using $f(x/n)=f(x)/n$ this gives continuity at $0$, hence everywhere. A continuous additive function is $f(x)=cx$, because it is $\mathbb Q$-linear and continuous.

\problem{3}
\hint{1} Call a sequence good if it is monotone and satisfies the doubling condition with respect to its first term. Strengthen the statement with two parameters: if $|X|\ge\binom{a+b}{a}+1$, then $X$ contains an increasing good sequence of length $a+2$ or a decreasing good sequence of length $b+2$. The problem is the case $a=b=k-2$. Induct on $a+b$, using Pascal's rule $\binom{a+b}{a}=\binom{a+b-1}{a-1}+\binom{a+b-1}{a}$.
\hint{2} Let $m=\min X$, $M=\max X$, $c=(m+M)/2$, and split $X$ into $L=\{x\le c\}$ and $U=\{x>c\}$. By Pascal's rule, $|L|\ge\binom{a-1+b}{a-1}+1$ or $|U|\ge\binom{a+b-1}{a}+1$. Apply the induction hypothesis to that half with $(a-1,b)$ or $(a,b-1)$ respectively. The base cases $a=0$ or $b=0$ just ask for two points.
\hint{3} Key observation: an increasing good sequence $y_1<\dots<y_r$ inside $L$ can be extended by $M$. You need $M-y_1\ge 2(y_r-y_1)$, i.e.\ $M-y_r\ge y_r-y_1$, and indeed $M-y_r\ge M-c=c-m\ge y_r-y_1$. Symmetrically, a decreasing good sequence inside $U$ can be extended by $m$. In each case either the half already contains the required sequence, or it contains one that is one term short in the direction that can be extended.

\problem{4}
\hint{1} Divisibility: the boxes are the odd parts. Write each number as $2^{s}u$ with $u$ odd. In $\{1,\dots,2n\}$ there are only $n$ odd values of $u$, and two numbers with the same odd part divide one another.
\hint{2} Monotone subsequences: the boxes are positions in an ordering. Give each term $a_i$ the label $\ell_i$, the length of the longest increasing subsequence ending at $a_i$. If some $\ell_i\ge n+1$ you are done. Otherwise $n^2+1$ terms have labels in $\{1,\dots,n\}$, so $n+1$ terms share a label.
\hint{3} If $i<j$ and $\ell_i=\ell_j$, then $a_i>a_j$, since otherwise $a_j$ would extend the sequence ending at $a_i$ and $\ell_j>\ell_i$. Hence the $n+1$ terms with equal label form a decreasing subsequence.

\problem{5}
\hint{1} Count pairs (colouring, monochromatic $k$-set). There are $2^{\binom n2}$ colourings in total. A fixed $k$-set of vertices is monochromatic in exactly $2\cdot 2^{\binom n2-\binom k2}$ of them. Equivalently: colour at random, and the expected number of monochromatic $K_k$ is $\binom nk 2^{1-\binom k2}$.
\hint{2} If this number is $<1$, the bad colourings cannot cover all colourings, so a good one exists. Check: $\binom nk 2^{1-\binom k2}<\frac{n^k}{k!}2^{1-k(k-1)/2}<\frac{2^{1+k/2}}{k!}\le 1$ for $k\ge 3$, using $n^k<2^{k^2/2}$.
\hint{3} The general principle: if the bad events $B_i$ satisfy $\sum_i P(B_i)<1$ (the union bound), some outcome avoids all of them. Equivalently, if (number of bad configurations) times (the largest number of objects each one spoils) is less than the total, a good object exists.

\end{document}
